\section{Introduction}
In the interior of the most massive of stars, three $\alpha$ particles are constantly fused together to become carbon-12. The energy released by this reaction counteracts the endless pressure from the colossal weight of the star's outer layers. Eventually gravity wins the battle, and the star meets a violent end, scattering the carbon into space. This process has seeded so much carbon throughout the Universe that it has become the fourth most abundant element by mass, with \ce{^12C} its most common form. With so much carbon available, it has become an essential building block of all known life; so essential that we can hardly imagine life without it. \\

Since some of the first research of the isotopes were carried out at the Cavendish Laboratory \note{ref}, it has been the subject of countless other studies. From the prediction and following discovery of the Hoyle state \note{ref}, so essential for stellar nucleosynthesis, to current medicinal studies of its uses in proton therapy \note{ref}, and to more experimental studies of its higher excited states; its importance can be clearly seen \note{sounds weird}. The goal of this thesis is to expand on the latter group, or more specifically on the understanding of the \SI{17.76}{MeV} $0^+$ state. Much work has already been done on the lower-lying \SI{16.62}{MeV} $2^-$ state \note{ref}, primarily in determining the branching ratios to \ce{^8Be}, but due to \note{equipment issues?}, this $0^+$ state has been somewhat neglected. 

\subsection{Brief review of the \SI{17.76}{MeV} state}
In the sequential picture, the \ce{^11B(p, 3\alpha)} reaction can proceed through both the ground state \ce{^11B(p, \alpha_0)^8Be} and the first excited state \ce{^11B(p, \alpha_1)^8Be^*} of \ce{^8Be}. One of the earliest measurements of the \SI{17.76}{MeV} resonance of \ce{^12C} were performed by Segel et al. \cite{Segel} back in 1965, where they noted that due to emission through the $\alpha_0$ channel, the state must have natural parity ($\pi = (-1)^J$). Based on a Legendre series fit to their data, they could also tentatively label it as a $0^+$ state, although they also saw a substantial $a_3$ component, indicating interference from the nearby \SI{18.35}{MeV} $3^-$ resonance. Later studies \cite{Hanna}\cite{Munch} supports this classification. \note{ask Hans about importance of e.g. Giorni, Bronson}. \\

Another relevant nearby state is the \SI{16.57}{MeV} $2-$ resonance, which was extensively studied by e.g. \cite{Kuhlwein}. 